\begin{proof}[Proof of \cref{obj:thm:finite-honest-upper}]
We first establish the bounds for \(n\ge2\). Fix an admissible arithmetic engine \(\mathsf E\), an admissible naming policy \(\mathsf N\), and public exponents satisfying \(0<\beta<1/4\) and \(\beta<\alpha\le1\).

\begin{enumerate}
\item Define the retained ranks, smoothness sum, and rate exponents by
\[
\begin{gathered}
k_C=k_C(n,\alpha,\beta;\mathsf E,\mathsf N),
\qquad
k_S=k_S(n,\beta;\mathsf E,\mathsf N),
\qquad s=\alpha+\beta,\\
\mathsf a(\alpha,\beta)
=\min\left\{\frac12,\frac{2s}{2s+1}\right\},
\qquad
\mathsf b(\beta)=\frac{4\beta}{4\beta+1}.
\end{gathered}
\]
These are the ranks and exponents of
\cref{obj:def:resolutions,obj:def:diagnostic-envelope}.
For \(k\ge1\), put
\[
V_n(k)=\frac{10}{n}+\frac{4k}{n(n-1)},
\qquad
b_C=8k_C^{-s}+\sqrt{20V_n(k_C)},
\qquad
b_S=25k_S^{-2\beta}+\sqrt{20V_n(k_S)}.
\]
Here \(b_C,b_S\) are the coordinate error radii, as in
\cref{obj:def:upper-handle}.

By \cref{obj:lem:fixed-budget-realization}, with target resolutions
\[
R_C=n^{2/(2s+1)},
\qquad R_S=n^{2/(4\beta+1)},
\]
the ranks satisfy
\[
1\le R_C\le k_C\le3R_C,
\qquad
1\le R_S\le k_S\le3R_S.
\]
In particular, both ranks are positive and are fixed before sampling.

For either pair
\[
(\gamma_{\mathrm{bal}},k_{\mathrm{bal}})
\in\{(s,k_C),(2\beta,k_S)\},
\]
the lower rank bound and \(\gamma_{\mathrm{bal}}>0\) give
\[
k_{\mathrm{bal}}^{-\gamma_{\mathrm{bal}}}
\le
n^{-2\gamma_{\mathrm{bal}}/(2\gamma_{\mathrm{bal}}+1)}.
\]
Also \(n(n-1)\ge n^2/2>0\), so the upper rank bound yields
\[
20V_n(k_{\mathrm{bal}})
\le
\frac{200}{n}
+
480\,\frac{n^{2/(2\gamma_{\mathrm{bal}}+1)}}{n^2}
=
200(n^{-1/2})^2
+
480\left(
n^{-2\gamma_{\mathrm{bal}}/(2\gamma_{\mathrm{bal}}+1)}
\right)^2.
\]
Both powers are nonnegative. The last expression is bounded by
\[
\left[
40\left(
n^{-1/2}
+
n^{-2\gamma_{\mathrm{bal}}/(2\gamma_{\mathrm{bal}}+1)}
\right)
\right]^2,
\]
because expansion of this square gives coefficients \(1600\) on both squared terms and a nonnegative cross term. Taking square roots therefore gives
\[
\sqrt{20V_n(k_{\mathrm{bal}})}
\le
40\left(
n^{-1/2}
+
n^{-2\gamma_{\mathrm{bal}}/(2\gamma_{\mathrm{bal}}+1)}
\right).
\]
Each power on the right is at most the power with exponent
\(-\min\{1/2,2\gamma_{\mathrm{bal}}/(2\gamma_{\mathrm{bal}}+1)\}\).
Since \(0<\beta<1/4\),
\[
\frac{4\beta}{4\beta+1}<\frac12.
\]
Consequently,
\[
\begin{aligned}
b_C
&\le88\,n^{-\mathsf a(\alpha,\beta)}
\le128\,n^{-\mathsf a(\alpha,\beta)},\\
b_S
&\le105\,n^{-\mathsf b(\beta)}
\le128\,n^{-\mathsf b(\beta)}.
\end{aligned}
\]
The exponents are positive, and \(\mathsf a(\alpha,\beta)\le1/2\le1\). Thus
\[
n^{-\mathsf a(\alpha,\beta)}\le1,
\qquad
n^{-\mathsf b(\beta)}\le1,
\qquad
\frac1n\le n^{-\mathsf a(\alpha,\beta)}.
\]
% lean: CausalSmith.Stat.LogoddsLowsmoothFrontier.coordinate_radii_power_bounds
% lean: CausalSmith.Stat.LogoddsLowsmoothFrontier.diagnostic_power_bounds

\item With the denominator floor and sensitivity coefficients defined by
\[
s_0=\frac1{512},
\qquad
\mathsf L_{\mathrm{sens}}=\frac{2\exp(1/2)}{s_0},
\qquad
\mathsf H_{\mathrm{sens}}=\frac{2\exp(1/2)}{s_0^2},
\]
choose
\[
K_0
=
128\mathsf L_{\mathrm{sens}}
+32768\mathsf H_{\mathrm{sens}}
+128\mathsf H_{\mathrm{sens}}\exp(1/2)+2.
\]
All its summands are nonnegative, so \(K_0\ge1>0\).
For every \(r\in[0,1/2]\), the preceding bounds imply
\[
\begin{aligned}
&\mathsf L_{\mathrm{sens}}b_C
+\mathsf H_{\mathrm{sens}}b_S
   \{\exp(1/2)r+2b_C\}
+\frac2n\\
&\quad\le
128\mathsf L_{\mathrm{sens}}n^{-\mathsf a(\alpha,\beta)}
+128\mathsf H_{\mathrm{sens}}n^{-\mathsf b(\beta)}
 \{\exp(1/2)r+256n^{-\mathsf a(\alpha,\beta)}\}
+2n^{-\mathsf a(\alpha,\beta)}\\
&\quad\le
(128\mathsf L_{\mathrm{sens}}
 +32768\mathsf H_{\mathrm{sens}}+2)
 n^{-\mathsf a(\alpha,\beta)}
+
128\mathsf H_{\mathrm{sens}}\exp(1/2)\,
r n^{-\mathsf b(\beta)}\\
&\quad\le
K_0\left\{
n^{-\mathsf a(\alpha,\beta)}
+r n^{-\mathsf b(\beta)}
\right\}
=
K_0d_n(\alpha,\beta;r).
\end{aligned}
\]
The second inequality uses
\(n^{-\mathsf b(\beta)}n^{-\mathsf a(\alpha,\beta)}
\le n^{-\mathsf a(\alpha,\beta)}\).
The last inequality follows by adding the nonnegative terms that complete the common coefficient \(K_0\).
This choice depends solely on the fixed denominator floor.
% lean: CausalSmith.Stat.LogoddsLowsmoothFrontier.radius_sensitivity_rate_assembly
% lean: CausalSmith.Stat.LogoddsLowsmoothFrontier.upper_radii_uniform_rate

\item Fix \(P\in\mathcal M(\alpha,\beta)\). Write
\[
\mathbb P_{P,n}=P^{\otimes n}\otimes\lambda,
\qquad
\mathbb E_{P,n}[\cdot]=\int(\cdot)\,d\mathbb P_{P,n},
\]
where \(\lambda\) is the uniform law of the independent seed. Variances below are under this same probability law.

We first control the unconditional absolute mean of the numerator estimate. Define, for \(x\in[0,1]\),
\[
\begin{gathered}
e(P,x)=\ell\{g(P,x)\},\\
\mu_0(P,x)=\ell\{\nu(P,x)\},
\qquad
\mu_1(P,x)=\ell\{\nu(P,x)+\theta(P)\},\\
m(P,x)=E_P[Y\mid X=x]
=(1-e(P,x))\mu_0(P,x)+e(P,x)\mu_1(P,x).
\end{gathered}
\]
The homogeneous relation and native-logit bounds are those of
\cref{obj:def:model}. The logistic derivative satisfies
\[
0<\ell'(t_{\ell})
=\ell(t_{\ell})\{1-\ell(t_{\ell})\}
\le\frac14
\qquad(t_{\ell}\in\mathbb R).
\]
The derivative bound on the interval between any two real arguments gives
\[
|\ell(t_{\ell})-\ell(t_{\ell}')|
\le\frac14|t_{\ell}-t_{\ell}'|
\qquad(t_{\ell},t_{\ell}'\in\mathbb R).
\]
Hence, for \(x,z\in[0,1]\),
\[
\begin{aligned}
|e(P,x)-e(P,z)|
&\le\frac12|x-z|^\alpha
\le\frac12|x-z|^\beta,\\
|\mu_0(P,x)-\mu_0(P,z)|
&\le\frac12|x-z|^\beta,\\
|\mu_1(P,x)-\mu_1(P,z)|
&\le\frac12|x-z|^\beta,\\
|\mu_1(P,z)-\mu_0(P,z)|
&\le\frac14|\theta(P)|\le\frac18.
\end{aligned}
\]
The comparison of the two covariate powers uses
\(|x-z|\le1\) and \(\beta<\alpha\); it also holds at \(x=z\), since both exponents are positive.

The marginal mean has the decomposition
\[
\begin{aligned}
m(P,x)-m(P,z)
={}&(1-e(P,x))\{\mu_0(P,x)-\mu_0(P,z)\}\\
&+e(P,x)\{\mu_1(P,x)-\mu_1(P,z)\}\\
&+\{e(P,x)-e(P,z)\}
  \{\mu_1(P,z)-\mu_0(P,z)\}.
\end{aligned}
\]
Since \(0<e(P,x)<1\), the preceding moduli imply
\[
\begin{aligned}
|m(P,x)-m(P,z)|
&\le
\{1-e(P,x)\}\frac12|x-z|^\beta
+e(P,x)\frac12|x-z|^\beta
+\frac12|x-z|^\beta\frac18\\
&=\frac9{16}|x-z|^\beta.
\end{aligned}
\]
Thus
\[
[e(P,\cdot)]_\alpha\le\frac12,
\qquad
[m(P,\cdot)]_\beta\le\frac9{16}.
\]
These functions are continuous: the propensity and marginal mean are sums of continuous conditional cells.

Define the empirical treated-outcome mean by
\[
\overline{AY}_n=\frac1n\sum_{i=1}^n A_iY_i.
\]
The coordinate definition in \cref{obj:def:coordinate-estimates} gives
\[
\widehat C_n
=\overline{AY}_n-\mathbb U_{n,k_C}(A,Y),
\qquad
\mathbb E_{P,n}\overline{AY}_n=E_P[AY].
\]
Using the projection expectation and residual identities from
\cref{obj:lem:projection-control}, we obtain
\[
\mathbb E_{P,n}\widehat C_n-C(P)
=
\left\langle
e(P,\cdot)-\Pi_{k_C}e(P,\cdot),
m(P,\cdot)-\Pi_{k_C}m(P,\cdot)
\right\rangle_{L^2(\lambda)}.
\]
Cauchy--Schwarz and the approximation bound in the same cited statement yield
\[
\begin{aligned}
|\mathbb E_{P,n}\widehat C_n-C(P)|
&\le
\|e(P,\cdot)-\Pi_{k_C}e(P,\cdot)\|_2
\|m(P,\cdot)-\Pi_{k_C}m(P,\cdot)\|_2\\
&\le
25\left(\frac12\right)\left(\frac9{16}\right)
k_C^{-(\alpha+\beta)}
=\frac{225}{32}k_C^{-s}
\le8k_C^{-s}.
\end{aligned}
\]
Here \(0<\alpha\le1\), \(0<\beta<1\), and \(k_C\ge1\), so the approximation bound applies throughout the stated range.

Independence of the records and \(0\le AY\le1\) give
\[
\operatorname{Var}_{P,n}(\overline{AY}_n)
=\frac{\operatorname{Var}_P(AY)}n
\le\frac1n.
\]
Also, \cref{obj:lem:projection-control} supplies
\[
\operatorname{Var}_{P,n}\{\mathbb U_{n,k_C}(A,Y)\}
\le\frac4n+\frac{2k_C}{n(n-1)}.
\]
The variance identities for a sum and a difference give
\[
\begin{aligned}
\operatorname{Var}_{P,n}(\widehat C_n)
={}&2\operatorname{Var}_{P,n}(\overline{AY}_n)
+2\operatorname{Var}_{P,n}\{\mathbb U_{n,k_C}(A,Y)\}\\
&-\operatorname{Var}_{P,n}
 \{\overline{AY}_n+\mathbb U_{n,k_C}(A,Y)\}\\
\le{}&\frac2n+2\left\{\frac4n+\frac{2k_C}{n(n-1)}\right\}
=V_n(k_C),
\end{aligned}
\]
where the inequality uses nonnegativity of the last variance.
The finite-rank statistics are square integrable, so all these identities and subsequent expectations are well defined.

Cauchy--Schwarz on the probability space gives
\[
\mathbb E_{P,n}
\left|\widehat C_n-\mathbb E_{P,n}\widehat C_n\right|
\le
\sqrt{\operatorname{Var}_{P,n}(\widehat C_n)}.
\]
Consequently,
\[
\begin{aligned}
\mathbb E_{P,n}|\widehat C_n|
&\le
|\mathbb E_{P,n}\widehat C_n|
+\sqrt{\operatorname{Var}_{P,n}(\widehat C_n)}\\
&\le
|C(P)|+8k_C^{-s}+\sqrt{V_n(k_C)}
\le |C(P)|+b_C.
\end{aligned}
\]
For the last step, \(V_n(k_C)>0\) and
\(\sqrt{V_n(k_C)}\le\sqrt{20V_n(k_C)}\).

By \cref{obj:lem:observable-odds},
\[
C(P)=\{\exp(\theta(P))-1\}S(P),
\qquad S(P)\ge s_0.
\]
The product defining \(S(P)\) is at most one pointwise, so
\[
S(P)
=\int p_{10}(P,x)p_{01}(P,x)\,\lambda(dx)\le1,
\]
where the two conditional off-diagonal probabilities are
\[
p_{10}(P,x)=P(A=1,Y=0\mid X=x),
\qquad
p_{01}(P,x)=P(A=0,Y=1\mid X=x).
\]
Since \(|\theta(P)|\le1/2\), the exponential derivative is bounded by
\(\exp(1/2)\) on the segment between \(0\) and \(\theta(P)\). Therefore
\[
|\exp(\theta(P))-1|
\le\exp(1/2)|\theta(P)|,
\qquad
|C(P)|\le\exp(1/2)|\theta(P)|.
\]
In particular, for \(P\in\mathcal M_r(\alpha,\beta)\),
\[
\mathbb E_{P,n}|\widehat C_n|
\le\exp(1/2)r+b_C.
\]
% lean: CausalSmith.Stat.LogoddsLowsmoothFrontier.cHat_integral_abs_le
% lean: CausalSmith.Stat.LogoddsLowsmoothFrontier.covarianceCoordinate_abs_le_effect

\item We next bound the interval length on every sample. Use the confidence endpoints and clipped inversion of \cref{obj:def:upper-handle}:
\[
\begin{gathered}
c_\pm=\widehat C_n\pm b_C,
\qquad
s_\pm=\operatorname{clip}_{[s_0,1]}(\widehat S_n\pm b_S),\\
F(c_{\mathrm{rect}},s_{\mathrm{rect}})
=
\log\!\left(
\operatorname{clip}_{[\exp(-1/2),\exp(1/2)]}
\left(1+\frac{c_{\mathrm{rect}}}{s_{\mathrm{rect}}}\right)
\right)
\quad
(c_{\mathrm{rect}}\in\mathbb R,\ s_{\mathrm{rect}}>0),\\
L=\min\{F(c_-,s_-),F(c_-,s_+)\},
\qquad
R=\max\{F(c_+,s_-),F(c_+,s_+)\}.
\end{gathered}
\]
Clipping is a composition of minimum and maximum with fixed constants, each of which contracts absolute differences. Thus
\[
s_0\le s_-,s_+,
\qquad
|s_--s_+|
\le|(\widehat S_n-b_S)-(\widehat S_n+b_S)|
=2b_S.
\]

For real numbers satisfying
\(\exp(-1/2)\le v_{\log}\le w_{\log}\), the inequality
\(\log t\le t-1\) for \(t>0\) gives
\[
\log w_{\log}-\log v_{\log}
=
\log\frac{w_{\log}}{v_{\log}}
\le
\frac{w_{\log}-v_{\log}}{v_{\log}}
\le
\exp(1/2)(w_{\log}-v_{\log}).
\]
Interchanging the two arguments handles the reverse order. Applying this bound to the clipped arguments, and using contraction of the exponential-range clipping, gives
\[
\left|
F(c_{\mathrm{rect}},s_{\mathrm{rect}})
-
F(c_{\mathrm{rect}}',s_{\mathrm{rect}}')
\right|
\le
\exp(1/2)
\left|
\frac{c_{\mathrm{rect}}}{s_{\mathrm{rect}}}
-
\frac{c_{\mathrm{rect}}'}{s_{\mathrm{rect}}'}
\right|
\]
for all real numerators and positive denominators.

For the rectangle points with
\[
c_{\mathrm{rect}},c_{\mathrm{rect}}'\in[c_-,c_+],
\qquad
s_{\mathrm{rect}},s_{\mathrm{rect}}'\in\{s_-,s_+\},
\]
we have
\[
|c_{\mathrm{rect}}-c_{\mathrm{rect}}'|\le2b_C,
\qquad
|c_{\mathrm{rect}}'|\le|\widehat C_n|+b_C,
\qquad
|s_{\mathrm{rect}}-s_{\mathrm{rect}}'|\le2b_S.
\]
Inserting the cross quotient with numerator
\(c_{\mathrm{rect}}'\) and denominator \(s_{\mathrm{rect}}\) yields
\[
\begin{aligned}
\left|
\frac{c_{\mathrm{rect}}}{s_{\mathrm{rect}}}
-
\frac{c_{\mathrm{rect}}'}{s_{\mathrm{rect}}'}
\right|
&\le
\frac{|c_{\mathrm{rect}}-c_{\mathrm{rect}}'|}
     {s_{\mathrm{rect}}}
+
\frac{|c_{\mathrm{rect}}'|
      |s_{\mathrm{rect}}-s_{\mathrm{rect}}'|}
     {s_{\mathrm{rect}}s_{\mathrm{rect}}'}\\
&\le
\frac{2b_C}{s_0}
+
\frac{2b_S(|\widehat C_n|+b_C)}{s_0^2}.
\end{aligned}
\]
It follows that every pairing of an upper-numerator corner and a lower-numerator corner satisfies the same difference bound. Choosing the corners attaining the maximum \(R\) and minimum \(L\), including either choice in a tie, gives
\[
R-L
\le
\mathsf L_{\mathrm{sens}}b_C
+
\mathsf H_{\mathrm{sens}}b_S(|\widehat C_n|+b_C).
\]
By \cref{obj:lem:fixed-budget-realization}, the literal rational output has additional length at most \(2/n\). Hence, pointwise,
\[
\left|
I^{\mathrm{up}}_{n,\mathsf E,\mathsf N}(\alpha,\beta)(\mathcal O,U)
\right|
\le
\mathsf L_{\mathrm{sens}}b_C
+
\mathsf H_{\mathrm{sens}}b_S(|\widehat C_n|+b_C)
+\frac2n.
\]
% lean: CausalSmith.Stat.LogoddsLowsmoothFrontier.inversion_four_corner_width_le
% lean: CausalSmith.Stat.LogoddsLowsmoothFrontier.upperOutput_length_le

\item For \(P\in\mathcal M_r(\alpha,\beta)\), integrate the pointwise bound. The ranks and radii are fixed before sampling, and the numerator estimate is integrable. Therefore
\[
\begin{aligned}
&\mathbb E_{P,n}
\left|
I^{\mathrm{up}}_{n,\mathsf E,\mathsf N}(\alpha,\beta)(\mathcal O,U)
\right|\\
&\quad\le
\mathsf L_{\mathrm{sens}}b_C
+
\mathsf H_{\mathrm{sens}}b_S
 \{\mathbb E_{P,n}|\widehat C_n|+b_C\}
+\frac2n\\
&\quad\le
\mathsf L_{\mathrm{sens}}b_C
+
\mathsf H_{\mathrm{sens}}b_S
 \{\exp(1/2)r+2b_C\}
+\frac2n\\
&\quad\le K_0d_n(\alpha,\beta;r).
\end{aligned}
\]
The last inequality is the deterministic calibration established above.
% lean: CausalSmith.Stat.LogoddsLowsmoothFrontier.upperInterval_expectedLength_le_radii

\item To establish honesty, take any \(P\in\mathcal M(\alpha,\beta)\).
The numerator bias and variance bounds have already been established. For the denominator, the conditional probabilities defined above satisfy
\[
p_{10}(P,x)=e(P,x)\{1-\mu_1(P,x)\},
\qquad
p_{01}(P,x)=\{1-e(P,x)\}\mu_0(P,x).
\]
All their factors lie in \([0,1]\). Adding and subtracting the product with the first factor evaluated at \(x\) and the second at \(z\) gives
\[
\begin{aligned}
|p_{10}(P,x)-p_{10}(P,z)|
&\le |e(P,x)-e(P,z)|
   +|\mu_1(P,x)-\mu_1(P,z)|\\
&\le |x-z|^\beta,\\
|p_{01}(P,x)-p_{01}(P,z)|
&\le |e(P,x)-e(P,z)|
   +|\mu_0(P,x)-\mu_0(P,z)|\\
&\le |x-z|^\beta.
\end{aligned}
\]
Thus both continuous functions have Hölder seminorm at exponent \(\beta\) bounded by one.

The denominator marks \(A(1-Y)\) and \((1-A)Y\) have conditional means \(p_{10}(P,\cdot)\) and \(p_{01}(P,\cdot)\).
The expectation and residual identities in \cref{obj:lem:projection-control} therefore give
\[
\mathbb E_{P,n}\widehat S_n-S(P)
=
-\left\langle
p_{10}(P,\cdot)-\Pi_{k_S}p_{10}(P,\cdot),
p_{01}(P,\cdot)-\Pi_{k_S}p_{01}(P,\cdot)
\right\rangle_{L^2(\lambda)}.
\]
Cauchy--Schwarz and the cited approximation bound imply
\[
\begin{aligned}
|\mathbb E_{P,n}\widehat S_n-S(P)|
&\le
\|p_{10}(P,\cdot)-\Pi_{k_S}p_{10}(P,\cdot)\|_2
\|p_{01}(P,\cdot)-\Pi_{k_S}p_{01}(P,\cdot)\|_2\\
&\le25k_S^{-2\beta}.
\end{aligned}
\]
The variance bound in the same cited statement applies to these bounded marks:
\[
\operatorname{Var}_{P,n}(\widehat S_n)
\le\frac4n+\frac{2k_S}{n(n-1)}
\le V_n(k_S).
\]
Together, the four coordinate moment bounds are
\[
\begin{gathered}
|\mathbb E_{P,n}\widehat C_n-C(P)|\le8k_C^{-s},
\qquad
|\mathbb E_{P,n}\widehat S_n-S(P)|\le25k_S^{-2\beta},\\
\operatorname{Var}_{P,n}(\widehat C_n)\le V_n(k_C),
\qquad
\operatorname{Var}_{P,n}(\widehat S_n)\le V_n(k_S).
\end{gathered}
\]
% lean: CausalSmith.Stat.LogoddsLowsmoothFrontier.coordinateEstimates_bias_identities
% lean: CausalSmith.Stat.LogoddsLowsmoothFrontier.numerator_projection_bias_bound
% lean: CausalSmith.Stat.LogoddsLowsmoothFrontier.denominator_projection_bias_bound
% lean: CausalSmith.Stat.LogoddsLowsmoothFrontier.coordinateEstimates_variance_bounds

\item Define the coordinate failure events by
\[
\begin{aligned}
\mathcal B_C
&=\{(\mathcal O,U):|\widehat C_n-C(P)|>b_C\},\\
\mathcal B_S
&=\{(\mathcal O,U):|\widehat S_n-S(P)|>b_S\}.
\end{aligned}
\]
The triangle inequality and the bias bounds imply
\[
\begin{aligned}
\mathcal B_C
&\subseteq
\left\{
|\widehat C_n-\mathbb E_{P,n}\widehat C_n|
\ge\sqrt{20V_n(k_C)}
\right\},\\
\mathcal B_S
&\subseteq
\left\{
|\widehat S_n-\mathbb E_{P,n}\widehat S_n|
\ge\sqrt{20V_n(k_S)}
\right\}.
\end{aligned}
\]
Both variance radii are strictly positive. Integrating each squared centered coordinate over its displayed event gives
\[
\begin{aligned}
20V_n(k_C)\,
\mathbb P_{P,n}
\left\{
|\widehat C_n-\mathbb E_{P,n}\widehat C_n|
\ge\sqrt{20V_n(k_C)}
\right\}
&\le\operatorname{Var}_{P,n}(\widehat C_n),\\
20V_n(k_S)\,
\mathbb P_{P,n}
\left\{
|\widehat S_n-\mathbb E_{P,n}\widehat S_n|
\ge\sqrt{20V_n(k_S)}
\right\}
&\le\operatorname{Var}_{P,n}(\widehat S_n).
\end{aligned}
\]
Consequently,
\[
\mathbb P_{P,n}(\mathcal B_C)\le\frac1{20},
\qquad
\mathbb P_{P,n}(\mathcal B_S)\le\frac1{20},
\qquad
\mathbb P_{P,n}(\mathcal B_C\cup\mathcal B_S)\le\frac1{10}.
\]

Outside their union,
\[
C(P)\in[c_-,c_+],
\qquad
\widehat S_n-b_S\le S(P)\le\widehat S_n+b_S.
\]
We have already shown \(s_0\le S(P)\le1\).
Clipping is nondecreasing and fixes every point of \([s_0,1]\), so
\[
s_-\le S(P)\le s_+.
\]
For a fixed positive denominator, \(F\) is nondecreasing in its numerator.
For a fixed real numerator \(c_{\mathrm{rect}}\), its behavior as the positive denominator varies follows from
\[
\begin{cases}
\displaystyle
\frac{c_{\mathrm{rect}}}{s_+}
\le\frac{c_{\mathrm{rect}}}{s_{\mathrm{rect}}}
\le\frac{c_{\mathrm{rect}}}{s_-},
& c_{\mathrm{rect}}\ge0,\\[6pt]
\displaystyle
\frac{c_{\mathrm{rect}}}{s_-}
\le\frac{c_{\mathrm{rect}}}{s_{\mathrm{rect}}}
\le\frac{c_{\mathrm{rect}}}{s_+},
& c_{\mathrm{rect}}<0,
\end{cases}
\qquad
s_{\mathrm{rect}}\in[s_-,s_+].
\]
The clipping and logarithm preserve these inequalities. Applying them first to \(c_-\), then to \(c_+\), and using numerator monotonicity yields
\[
\begin{aligned}
L
&\le F(c_-,S(P))
\le F(C(P),S(P))\\
&\le F(c_+,S(P))
\le R.
\end{aligned}
\]
By \cref{obj:lem:observable-odds} and the effect envelope,
\[
1+\frac{C(P)}{S(P)}
=\exp(\theta(P))
\in[\exp(-1/2),\exp(1/2)],
\qquad
F(C(P),S(P))=\theta(P).
\]
Thus \(\theta(P)\in[L,R]\) outside \(\mathcal B_C\cup\mathcal B_S\).
The containment supplied by \cref{obj:lem:fixed-budget-realization} gives
\[
[L,R]\subseteq
I^{\mathrm{up}}_{n,\mathsf E,\mathsf N}(\alpha,\beta)(\mathcal O,U)
\]
on every sample. Hence
\[
\begin{aligned}
1
&\le
\mathbb P_{P,n}
\left\{
\theta(P)\in
I^{\mathrm{up}}_{n,\mathsf E,\mathsf N}(\alpha,\beta)(\mathcal O,U)
\right\}
+\mathbb P_{P,n}(\mathcal B_C\cup\mathcal B_S)\\
&\le
\mathbb P_{P,n}
\left\{
\theta(P)\in
I^{\mathrm{up}}_{n,\mathsf E,\mathsf N}(\alpha,\beta)(\mathcal O,U)
\right\}
+\frac1{10}.
\end{aligned}
\]
The coverage probability is therefore at least \(9/10\).
The Borel output property from the same cited statement gives measurable coverage events and interval length. Since \(P\) was arbitrary in the ambient model, this proves
\[
I^{\mathrm{up}}_{n,\mathsf E,\mathsf N}(\alpha,\beta)
\in\mathcal A_n(\alpha,\beta).
\]
% lean: CausalSmith.Stat.LogoddsLowsmoothFrontier.coordinate_error_probability_le
% lean: CausalSmith.Stat.LogoddsLowsmoothFrontier.upperInterval_coverage_of_coordinate_moments

\item Now let \(n\ge1\). If \(n=1\), the first branch of
\cref{obj:def:upper-handle} returns
\[
I^{\mathrm{up}}_{1,\mathsf E,\mathsf N}(\alpha,\beta)(\mathcal O,U)
=[-1/2,1/2]
\]
at every input. The effect envelope gives coverage one, and its length is one. For every \(r\in[0,1/2]\),
\[
d_1(\alpha,\beta;r)=1+r,
\qquad
\mathbb E_{P,1}
\left|
I^{\mathrm{up}}_{1,\mathsf E,\mathsf N}(\alpha,\beta)(\mathcal O,U)
\right|
=1
\le K_0(1+r).
\]
If \(n\ne1\), the integer assumption \(n\ge1\) implies \(n\ge2\), and the honesty and expected-length bounds established above apply. Thus the same positive constant \(K_0\) gives both assertions for every admissible engine, policy, exponent pair, sample size, radius, and law in the theorem.
% lean: CausalSmith.Stat.LogoddsLowsmoothFrontier.upperInterval_one_honest
% lean: CausalSmith.Stat.LogoddsLowsmoothFrontier.upperInterval_one_expectedLength_bound

\item Finally, every branch of the literal output construction either takes a nonempty intersection of a rational endpoint box with the full effect box or returns the full effect box. Each resulting box has ordered rational endpoints and is contained in \([-1/2,1/2]\); the output includes both endpoints. Consequently, for every sample and seed there exist \(l,u\in\mathbb Q\) such that
\[
I^{\mathrm{up}}_{n,\mathsf E,\mathsf N}(\alpha,\beta)(\mathcal O,U)
=[l,u],
\qquad
-\frac12\le l\le u\le\frac12.
\]
For \(n=1\), the full-box branch gives the asserted equality with
\([-1/2,1/2]\) at every sample and seed.
% lean: CausalSmith.Stat.LogoddsLowsmoothFrontier.upperOutput_rational_shape
% lean: CausalSmith.Stat.LogoddsLowsmoothFrontier.upperInterval_one_set
\end{enumerate}
\end{proof}
